\section{NLP 实验实践建议}

课程的后半部分，Manning 教授给出了关于 NLP 研究和实验的实用建议。

\subsection{研究项目的基本要素}

\begin{importantbox}{好的 NLP 研究项目的关键要素}
\begin{enumerate}
    \item \textbf{明确的数据来源}：在项目开始前就确定使用什么数据
    \item \textbf{清晰的评估方法}：定义如何衡量系统性能
    \item \textbf{合适的基线}（Baseline）：必须有一个对比对象来证明你的方法有价值
    \item \textbf{实质性的价值贡献}：不能只是下载一个好模型跑一下数据——需要有分析、理解、改进
\end{enumerate}
\end{importantbox}

\subsection{基线的重要性}

Manning 教授强调，任何研究都需要一个合适的基线：

\begin{itemize}
    \item 如果之前有人做过相同任务，使用他们的结果作为基线
    \item 如果是全新任务，设计一个简单直觉的方法作为基线（例如：用词向量平均值的余弦相似度做文本相似度）
    \item 你的复杂模型必须\textbf{显著优于}这个简单基线，否则不具说服力
\end{itemize}

\begin{warningbox}{没有基线的研究是不完整的}
仅仅展示``我的模型达到了某个数字''是不够的。没有对比，读者无法判断这个数字是好还是坏。Manning 教授举例：如果你构建了一个复杂的神经网络来计算文本相似度，但它的表现还不如简单地将词向量取平均后做点积，那么这个系统就不是一个好系统。
\end{warningbox}

\subsection{计算资源的策略}

在 GPU 资源日益紧张的今天（Manning 教授戏称``这都怪 OpenAI''），合理利用计算资源至关重要：

\begin{itemize}
    \item \textbf{云端 Notebook}：Google Colab、Kaggle Notebooks、AWS SageMaker Studio Lab 提供免费 GPU
    \item \textbf{低成本 GPU 提供商}：Modal、Vast.ai 等
    \item \textbf{API 访问}：对于 LLM 项目，API 调用（如 Together AI）可能比自己训练更高效
    \item \textbf{选择合适的模型大小}：7B 参数模型的 API 成本远低于 70B 模型，在许多任务上已经够用
\end{itemize}

\begin{knowledgebox}{模型大小的权衡}
Manning 教授建议：在做实验之前，先考虑你\textbf{真正需要}多大的模型。如果 7B 参数模型足以证明你的观点，就不必使用更大的模型。50美元的 API 额度在 7B 模型上可以处理海量 token，但在大模型上可能很快用完。研究的关键不在于使用最大的模型，而在于\textbf{证明有趣的结论}。
\end{knowledgebox}

\subsection{本章小结}

\begin{itemize}
    \item 好的研究需要明确的数据、评估方法和基线
    \item 基线是不可或缺的——没有对比就没有说服力
    \item 合理选择模型大小和计算策略，可以在有限预算下完成高质量研究
    \item 批判性思维是研究的核心能力：理解方法的优缺点，而非简单复述
\end{itemize}

%% ========== 总结与延伸 ==========

